% Partie : sets-functions-relations
% Chapitre : relations
% Section : reflections
%
\documentclass[../../../include/open-logic-section]{subfiles}

\begin{document}

\olfileid[fr]{sfr}{rel}{ref}
\olsection{Réflexions philosophiques}

Dans la \olref[set]{sec}, nous avons défini les relations comme
certains ensembles. Arrêtons-nous un instant sur une question
philosophique : que \emph{fait} une telle définition ?
Il est fort douteux que nous voulions dire avoir
\emph{découvert} des identités métaphysiques : par exemple,
que la relation d'ordre sur $\Nat$ \emph{se serait révélée être}
l'ensemble $R= \Setabs{\tuple{n,m}}{n, m \in \Nat\text{ et }
n < m}$ défini dans la \olref[set]{sec}.
Voici trois raisons d'en douter.

Premièrement, dans la \olref[set][pai]{wienerkuratowski},
nous avons défini $\tuple{a, b} = \{\{a\}, \{a, b\}\}$.
Considérons plutôt la définition
$\lVert a, b\rVert = \{\{b\}, \{a, b\}\} = \tuple{b,a}$.
Lorsque $a \neq b$, on a
$\tuple{a, b} \neq \lVert a,b\rVert$.
Pourtant, nous aurions tout aussi bien pu choisir
$\lVert a,b\rVert$ comme définition d'un couple, au lieu de
$\tuple{a,b}$. Les deux définitions auraient aussi bien convenu.
Nous avons donc deux candidats tout aussi satisfaisants
pour «~être~» la relation d'ordre sur les entiers naturels :
\begin{align*}
		R &= \Setabs{\tuple{n,m}}{n, m \in \Nat \text{ et }n < m}\\
		S &= \Setabs{\lVert n,m\rVert}{n, m \in \Nat \text{ et }n < m}.
\end{align*}
Puisque $R \neq S$, par extensionnalité, ils ne peuvent
manifestement pas être \emph{tous deux} identiques à la relation
d'ordre sur~$\Nat$. Mais il serait arbitraire, et donc quelque
peu embarrassant, d'affirmer que c'est $R$, plutôt que $S$
(ou l'inverse), qui \emph{est}, dans les faits, cette relation
d'ordre. Il s'agit d'un exemple très simple d'un argument
contre le réductionnisme ensembliste, rendu célèbre par
Benacerraf en \citeyear{Benacerraf1965}.
Nous y reviendrons à plusieurs reprises.

Deuxièmement, si \emph{toute} relation doit être identifiée
à un ensemble, la relation d'appartenance elle-même, $\in$,
devrait être un ensemble particulier : précisément,
$\Setabs{\tuple{x,y}}{x \in y}$.
Mais cet ensemble existe-t-il ?
Le paradoxe de Russell nous avertit que son existence
ne va pas de soi. En fait,
\oliflabeldef{cumul:::part}{la théorie des ensembles développée
dans la partie \olref[cumul][][]{part} \emph{exclura} l'existence
de cet ensemble.\footnote{En anticipant sur la suite, voici
pourquoi. Supposons, pour raisonner par l'absurde, que
$I = \Setabs{\tuple{x,y}}{x \in y}$ existe.
Alors $\bigcup \bigcup I$ est l'ensemble universel,
ce qui contredit le résultat \olref[sfr][z][sep]{thm:NoUniversalSet}.}}{on peut
développer rigoureusement la théorie des ensembles sous forme
axiomatique, et cette théorie exclura effectivement l'existence
de cet ensemble.}
Ainsi, même si certaines relations peuvent être traitées
comme des ensembles, la relation d'appartenance devra
constituer un cas particulier.

Troisièmement, en «~identifiant~» les relations à des ensembles,
nous nous sommes autorisés à écrire $Rxy$ pour
$\tuple{x,y} \in R$. Cela convient si la relation
d'appartenance, «~$\in$~», est traitée \emph{comme} un prédicat.
Mais si «~$\in$~» désigne un certain type d'ensemble,
l'expression «~$\tuple{x,y} \in R$~» ne consiste plus
qu'en trois termes singuliers désignant des ensembles :
«~$\tuple{x,y}$~», «~$\in$~» et «~$R$~».
Or une telle liste de noms n'est pas plus capable d'exprimer
une proposition que la suite de mots dépourvue de sens :
«~la tasse le pot à crayons la table~».
Là encore, même si certaines relations peuvent être traitées
comme des ensembles, la relation d'appartenance doit constituer
un cas particulier. Cet argument réunit une version simple
du paradoxe frégéen du concept \emph{cheval} et une célèbre
objection que Wittgenstein a adressée à Russell.

Que faut-il en conclure ?
Il n'y a rien de \emph{fautif} à dire que les relations sur
les nombres sont des ensembles, à condition de comprendre
dans quel esprit on le dit. Nous n'affirmons pas une identité
métaphysique. Nous remarquons simplement que, dans certains
contextes, on peut \emph{traiter} certaines relations comme
certains ensembles, et que c'est ce que nous ferons.

\end{document}
